% Kapitel 5 --- Methode
% Budget: 3 Seiten --- das Herzstück, hier entscheidet sich die Note
% Quelle im Vault: 20 Recherche/Methodik

\section{Methode}
\label{sec:methode}

\subsection{Analyseeinheit: der Sachverhalt}

\todo{Definition als Tupel aus Entität, Attribut, Wert, Gültigkeitsintervall, Quelle und
Fundstelle. Ein Beispiel ausgeschrieben. Begründung: Ohne diese Definition bleibt
\enquote{derselbe Sachverhalt} Alltagssprache und die Messung wird nicht prüfbar.}

\subsection{Extraktion mit Zitatpflicht}

\todo{Sprachmodell mit festem Schema, jeder Wert mit Belegstelle; Werte ohne Fundstelle
gelten als nicht extrahiert. Begründung: Ein Modell liefert plausible Werte auch dort,
wo keine stehen. Temperatur, Prompt-Version, Modellversion --- gehört in den Anhang.}

\subsection{Normalisierung}

\todo{Datum, Zahl, Einheit, Entitätsalias, Scope. Tabelle der Regeln mit je einem
Beispiel. Hinweis: Die Aliasliste ist selbst ein Ergebnis, weil sie zeigt, wie viele
Schreibweisen derselbe Studiengang hat.}

\subsection{Abweichungstypologie}

\todo{Die fünf Typen (echter Widerspruch, Versionsunterschied, Scope-Unterschied,
Formulierung, Extraktionsfehler) mit Entscheidungsregel je Typ. Das ist der eigene
Beitrag --- hier ausführlich werden.}

\subsection{Goldstandard und Reliabilität}

\todo{Stichprobe, zwei Durchgänge, Konfliktlösung, Krippendorffs Alpha mit Konvention
\parencite{krippendorff2004}. Und die Validierung des Modells gegen menschliche Labels
\parencite{bavaresco2025llms}.}

\subsection{Metriken}

\todo{Precision, Recall, F1 je Sachverhaltstyp; Wilson-Konfidenzintervall; Formeln
hinschreiben, nicht nur nennen. Warum je Typ und nicht global.}
